\setcounter{section}{0}
\section*{四种常见构型}
\phantomsection
\addcontentsline{toc}{section}{四种常见构型}
四图中的 \(ABCD\) 均为正方形，各动点在指定边的内部。识别时先找等长邻边和直角，再看 \(45^\circ\) 在顶点还是交点。原图与辅助线图并排，辅助线用虚线表示。

\block{图 A 顶点夹半角：完整证明}
\textbf{条件与结论}\quad \(E\in BC\)、\(F\in CD\)，\(\angle EAF=45^\circ\)。则 \(EF=BE+DF\)；反过来，\(EF=BE+DF\) 也能推出 \(\angle EAF=45^\circ\)。
\diagram{\figIntroA\hspace{1.2cm}\figIntroASolve}

\think 把 \(BE\) 搬到 \(CD\) 的延长线上，使 \(BE+DF\) 成为一条线段，再用半角证明这条线段等于 \(EF\)。
\solve 延长 \(CD\) 至 \(G\)，使 \(DG=BE\)，连接 \(AG\)。点序为 \(C,F,D,G\)。这相当于把 \(\triangle ABE\) 绕 \(A\) 逆时针旋转 \(90^\circ\) 到 \(\triangle ADG\)。
\begin{enumerate}
\item 在 \(\triangle ABE\) 与 \(\triangle ADG\) 中，\(AB=AD\)、\(BE=DG\)、\(\angle ABE=\angle ADG=90^\circ\)，所以 \(\triangle ABE\cong\triangle ADG\)（SAS）。于是 \(AE=AG\)，\(\angle BAE=\angle DAG\)。
\item 因为 \(\angle BAE+\angle EAD=90^\circ\)，所以
\[
\angle EAG=\angle EAD+\angle DAG=90^\circ.
\]
\(AF\) 在 \(\angle EAG\) 内，而 \(\angle EAF=45^\circ\)，故 \(\angle FAG=45^\circ\)。由 \(AE=AG\)、\(AF\) 公共、\(\angle EAF=\angle GAF\)，得 \(\triangle EAF\cong\triangle GAF\)（SAS），因此 \(EF=FG\)。
\item 由 \(F,D,G\) 的点序，\(EF=FG=FD+DG=DF+BE\)，结论成立。
\end{enumerate}
\insight{等长邻边提供第一次全等，半角提供第二次全等，最后按点序写线段和。}

\clearpage
\block{图 A 的逆命题与相应结论}
\textbf{逆命题的证明}\quad 仍作 \(DG=BE\)，第一组全等仍给出 \(AE=AG\)、\(\angle EAG=90^\circ\)。若 \(EF=BE+DF\)，则 \(EF=FG\)。结合 \(AF\) 公共，\(\triangle EAF\cong\triangle GAF\)（SSS），所以 \(\angle EAF=\angle GAF\)。两角相加为 \(90^\circ\)，故 \(\angle EAF=45^\circ\)。

设 \(AB=s\)、\(BE=e\)、\(DF=f\)，其中 \(0<e,f<s\)。由线段和可以继续得到：
\begin{enumerate}
\item \textbf{定周长。} \(CE=s-e\)、\(CF=s-f\)，所以
\[
EF=e+f\quad\Longrightarrow\quad
C_{\triangle CEF}=(s-e)+(s-f)+(e+f)=2s.
\]
\item \textbf{长度关系。} 在 \(\mathrm{Rt}\triangle CEF\) 中，由勾股定理，
\[
(s-e)^2+(s-f)^2=(e+f)^2.
\]
\item \textbf{已知 \(AB\)、\(BE\) 求另两段。} 由上式解得
\[
\boxed{DF=\frac{s(s-e)}{s+e}},\qquad
\boxed{CF=s-DF=\frac{2se}{s+e}}.
\]
\end{enumerate}

\par\bigskip
\noindent\begin{minipage}{\linewidth}
\block{图 B 交叉角：平行转移}
\(E\in BC\)、\(F\in AB\)，\(AE\) 与 \(CF\) 交于 \(P\)，\(\angle EPC=45^\circ\)。过 \(A\) 作 \(AG\parallel CF\)，交 \(CD\) 于 \(G\)，并连接 \(EG\)。
\diagram{\figIntroB\hspace{1.2cm}\figIntroBSolve}
\[
AG\parallel CF\quad\Longrightarrow\quad\angle EAG=\angle EPC=45^\circ.
\]
\[
\left.\begin{gathered}
AD=CB,\quad\angle ADG=\angle CBF=90^\circ\\
\angle DAG=\angle BCF
\end{gathered}\right\}
\ \Longrightarrow\ \triangle ADG\cong\triangle CBF\ (\mathrm{ASA})
\ \Longrightarrow\ DG=BF.
\]
\[
\text{转为图 A}\quad\Longrightarrow\quad
\boxed{EG=BE+DG=BE+BF}.
\]
若 \(AB=s\)、\(BE=e\)、\(BF=f\)，则由图 A 的勾股等式化简得 \(\boxed{ef+s(e+f)=s^2}\)。
\end{minipage}

\clearpage
\noindent\begin{minipage}{\linewidth}
\block{图 C 交叉角：垂直转移}
\(E\in BC\)、\(F\in AD\)，\(AE\) 与 \(BF\) 交于 \(P\)，\(\angle APF=45^\circ\)。过 \(A\) 作 \(AG\perp BF\)，交 \(CD\) 于 \(G\)，并连接 \(EG\)。
\diagram{\figIntroC\hspace{1.2cm}\figIntroCSolve}
\[
AG\perp BF,\quad\angle APF=45^\circ
\quad\Longrightarrow\quad\angle EAG=90^\circ-45^\circ=45^\circ.
\]
\[
\left.\begin{gathered}
AD=BA,\quad\angle ADG=\angle BAF=90^\circ\\
\angle DAG=\angle ABF
\end{gathered}\right\}
\ \Longrightarrow\ \triangle ADG\cong\triangle BAF\ (\mathrm{ASA})
\ \Longrightarrow\ DG=AF.
\]
其中 \(\angle DAG=\angle ABF\)，因为这两个锐角的两边分别垂直。
\[
\text{转为图 A}\quad\Longrightarrow\quad
\boxed{EG=BE+DG=BE+AF}.
\]
若 \(AB=s\)、\(BE=e\)、\(AF=f\)，则 \(\boxed{ef+s(e+f)=s^2}\)。
\end{minipage}

\par\bigskip
\noindent\begin{minipage}{\linewidth}
\block{图 D 两条线同时平移}
\(E\in AD\)、\(F\in BC\)、\(M\in AB\)、\(N\in CD\)，\(EF\) 与 \(MN\) 交于 \(P\)，\(\angle FPN=45^\circ\)。按图示位置，\(BF>AE\)、\(DN>AM\)。过 \(A\) 作 \(AT\parallel EF\)、\(AU\parallel MN\)，分别交 \(BC\)、\(CD\) 于 \(T,U\)，并连接 \(TU\)。
\diagram{\figIntroD\hspace{1.2cm}\figIntroDSolve}
\[
\begin{aligned}
AE\parallel TF,\ AT\parallel EF
&\ \Longrightarrow\ AEFT\text{ 为平行四边形}\\
&\ \Longrightarrow\ TF=AE\ \Longrightarrow\ BT=BF-AE;\\[5pt]
AM\parallel NU,\ AU\parallel MN
&\ \Longrightarrow\ AMNU\text{ 为平行四边形}\\
&\ \Longrightarrow\ NU=AM\ \Longrightarrow\ DU=DN-AM.
\end{aligned}
\]
\[
AT\parallel PF,\ AU\parallel PN
\quad\Longrightarrow\quad\angle TAU=45^\circ
\quad\Longrightarrow\quad\text{转为图 A}.
\]
\[
\boxed{TU=BT+DU=(BF-AE)+(DN-AM)}.
\]
若 \(AB=s\)、\(BT=e\)、\(DU=f\)，仍有 \(ef+s(e+f)=s^2\)。点序改变时，要重新判断射线方向及线段加减。
\end{minipage}
\clearpage
